% !TeX root = ../main.tex
\section{A Head Identity Extension}

% TODO: define equation for h
% TODO: define hysteresis function
% TODO: define coupling h -> rho
% TODO: run simulations showing stable heads away from Turing-selected positions
% TODO: test competition between nearby heads

\begin{itemize}
    \item Problem to solve:
    \begin{itemize}
        \item Dynamic source density can preserve memory for long times but may not create true stationary anchoring.
        \item Activator peaks are transient biochemical signals.
        \item A head should represent a more stable morphological commitment.
    \end{itemize}
    \item Proposed mechanism:
    \begin{itemize}
        \item Introduce a new non-diffusing variable \(h(t,x)\), called head identity.
        \item The activator \(u\) can induce \(h\), but only if \(u\) is high enough for long enough.
        \item \(h\) then generates or organizes source density \(\rho\).
    \end{itemize}
    \item Conceptual hierarchy:
    \begin{itemize}
        \item \(u\): fast activator signal;
        \item \(v\): inhibitor;
        \item \(h\): slow local head identity/morphological commitment;
        \item \(\rho\): positional field generated by heads.
    \end{itemize}
    \item Key modelling principle:
    \begin{itemize}
        \item Source density should encode distance from actual heads, not from temporary activator peaks.
    \end{itemize}
    \item Hysteresis:
    \begin{itemize}
        \item head formation threshold is higher than head maintenance threshold;
        \item it is harder to create a head than to maintain an existing head;
        \item after \(h\) is established, it does not disappear immediately when \(u\) decreases.
    \end{itemize}
    \item Possible implementation:
    \begin{itemize}
        \item \(h\) does not diffuse;
        \item \(h\) grows when \(u\) exceeds a high threshold;
        \item \(h\) decays when \(u\) falls below a lower threshold;
        \item use a smooth bistable or hysteretic response function.
    \end{itemize}
    \item Expected behavior:
    \begin{itemize}
        \item wound-induced activator peaks can create a head if strong and persistent enough;
        \item once a head is established, it anchors source density;
        \item peaks can remain at graft-induced positions rather than only Turing-selected positions;
        \item heads too close to each other may still compete and one may disappear.
    \end{itemize}
    \item Main claim:
    \begin{itemize}
        \item Separating activator peaks from head identity provides an autonomous mechanism for positional memory.
    \end{itemize}
\end{itemize}
